\documentclass[Screen16to9,17pt]{foils}
\usepackage{zencurity-slides}
\externaldocument{software-security-exercises}
\selectlanguage{english}

\begin{document}

\mytitlepage
{9. Program Building Blocks}
{EK  Software Security}


\slide{Goals for today}

\hlkimage{6cm}{thomas-galler-hZ3uF1-z2Qc-unsplash.jpg}

Todays goals:
\begin{list2}
\item Talk about the Design Patterns concept in a broader sense
\item Present some decisions concerning security and running applications securely
\item Patterns can also be found in other areas, like Patterns in Network Architecture
\end{list2}

  Photo by Thomas Galler on Unsplash

\slide{Plan for today}

\begin{list1}
\item Subjects
\begin{list2}
\item Common constructs
\item Recurring code patterns
\item Example programs with flaws: OpenSSH, OpenSSL, Windows MS-RPC DCOM, Linux teardrop
\end{list2}
\item Exercises
\begin{list2}
\item Static analysis with PMD
\item SBOM Scanning
\item Use a XML library in Python
\end{list2}
\end{list1}


\slide{Time schedule}

Plan for the day
\begin{list1}
\item 12:45 - 13:45 -- Session 0: Teaching: Code Patterns and Running Software Securely
\item 13:45 - 14:00 -- 15min break
\item 14:00 - 14:45 -- Session 1: Exercises: PMD Static Analysis
\item 14:45 - 15:00 -- 15min break
\item 15:00 - 15:45 -- Session 2: Teaching: Sample Infrastructure, API Versioning and Containers
\item 15:45 - 16:00 -- 15min break
\item 16:00 - 16:45 -- Session 3: Exercises: SBOM Scanning and XML Library in Python

\vskip 1em
\item 12:45 - 16:20 I will be here at EK
\end{list1}

\slide{Session 0: Teaching: Code Patterns and Running Software Securely}




\slide{Design Patterns}

\begin{list2}
\item \emph{Design Patterns: Elements of Reusable Object-Oriented Software} (1994), Erich Gamma, Richard Helm, Ralph Johnson, and John Vlissides

\item The book describes 23 classic software design patterns

\item \url{https://en.wikipedia.org/wiki/Design_Patterns}

\item Ideas of patterns precede this book, but became a more popular subject
\end{list2}



\slide{Common constructs}


\begin{list2}
\item Programs exhibit the same patterns, some examples:
\item Solve problems in the same domains
\item Need to store lists of strings / characters etc.
\item Data structures becomes useful in other programs
\item Sorting routines needed in many programs
\end{list2}





\slide{Auditing code patterns}

Various pitfalls, and areas
\begin{list2}
\item {\bf Variable Relationship}
\item Examples presented are C buffers
\item A variable-pointer to the buffer and the variable with the length
\item Their relationship can easily get invalidated, leading to vulnerabilities
\item Question: would object oriented programming help?
\end{list2}



\slide{Structure and Object Mismanagement}


\begin{list2}
\item Structures in C can help group related data elements
\item Both auditors and attackers can benefit ...
\item Object oriented programs and languages encapsulate this
\item Responsibility is on the object implementation, good!
\item We saw structures last time
\end{list2}



\slide{Uninitialized Variables}

\begin{minted}[fontsize=\footnotesize]{c}
    Packet *p = SCMalloc(SIZE_OF_PACKET);
    if (unlikely(p == NULL)) {
         return 0;
    }
    ThreadVars tv;
    DecodeThreadVars dtv;

    memset(&dtv, 0, sizeof(DecodeThreadVars));
    memset(&tv,  0, sizeof(ThreadVars));
    memset(p, 0, SIZE_OF_PACKET);
\end{minted}

\begin{list2}
\item Always make sure variables have well defined values.
\item Defensive program above use memset to clear contents of buffers
\item Example from Suricata allocating memory, checking it got the memory, clearing contents
\end{list2}


\slide{Manipulating Lists - and other data structures}


\begin{list2}
\item Storing data in a list is nice, can add and remove elements
\item Single linked list, start from head and go through list ... slow
\item Double linked list can go back again from element
\item More work updating, moving more pointers ... complex, may introduce errors
\item Example data structure double linked list\\
\url{https://algorithms.tutorialhorizon.com/doubly-linked-list-complete-implementation/}
\item Multiple problems can arise, also using the wrong structure for something can result in vulnerabilities
\end{list2}

\vskip 1cm
\centerline{Dont write your data structure libraries yourself}

\slide{Linux Teardrop}

\begin{minted}[fontsize=\footnotesize]{python}
#! /usr/bin/env python3
if attack == '0':
  print "Using attack 0"
  size=36    offset=3   load1="\x00"*size
  i=IP()   i.dst=target   i.flags="MF"   i.proto=17

  size=4     offset=18  load2="\x00"*size
  j=IP()   j.dst=target   j.flags=0   j.proto=17   j.frag=offset
  send(i/load1)
  send(j/load2)
\end{minted}

\begin{list2}
\item IP fragments, packets are split when crossing a link with lower MTU
\item If fragments created by an attacker are overlapping it creates a problem
\item Scapy example code by Sam Bowne, full code can be found at:\\
\url{https://samsclass.info/123/proj10/teardrop.htm}
\item what is a Maximum Transmission Unit (MTU)\\
See \url{https://en.wikipedia.org/wiki/Maximum_transmission_unit} for a description,\\ related/similar attack setting source=destination \url{https://en.wikipedia.org/wiki/LAND}
\end{list2}


\slide{Other problems}


\begin{list2}
\item Hashing algorithms
\item There have been multiple problems with hashing algorithms
\item Denial of Service and arbitrary code can be the result
\item Example vulns from popular programming languages, others have similar!\\
\url{https://www.cvedetails.com/vulnerability-list/vendor_id-74/product_id-128/cvssscoremin-7/cvssscoremax-7.99/PHP-PHP.html} search for hash \\
\url{https://github.com/bk2204/php-hash-dos} specific example in PHP 7.0.0~rc3-3\\
\url{https://www.ruby-lang.org/en/security/} Ruby has some
\item also when programmers selected wrong, or weak, hashing algorithms for passwords
\end{list2}

\slide{Control Flow}

\hlkimage{18cm}{ssl-tls-breaks-timeline.png}

\begin{list2}
\item If constructs, make sure to exercise each path
\item Case/switch constructs, make sure to catch a default
\item Loops being subverted to create buffer overflow, terminating conditions
\item Forgetting a \verb+break+ in a case
\item Picture selected due to the Apple "goto fail" - certificate validation\\
\link{https://dwheeler.com/essays/apple-goto-fail.html}
\end{list2}


\slide{Apple Goto Fail CVE-2014-1266}

\begin{minted}[fontsize=\footnotesize]{python}
if ((err = SSLHashSHA1.update(&hashCtx, &clientRandom)) != 0)
    goto fail;
if ((err = SSLHashSHA1.update(&hashCtx, &serverRandom)) != 0)
    goto fail;
if ((err = SSLHashSHA1.update(&hashCtx, &signedParams)) != 0)
    goto fail;
    goto fail;  /* MISTAKE THIS LINE SHOULD NOT BE HERE */
if ((err = SSLHashSHA1.final(&hashCtx, &hashOut)) != 0)
    goto fail;
\end{minted}

\begin{list2}
\item Always going to goto fail, second goto is always active
\item Leading to later checks not performed
\item Example code from\emph{Anatomy of a “goto fail” – Apple’s SSL bug explained}\\ {\footnotesize\url{https://nakedsecurity.sophos.com/2014/02/24/anatomy-of-a-goto-fail-apples-ssl-bug-explained-plus-an-unofficial-patch/}}

\end{list2}


\slide{Loop running over buffer}

\begin{list2}
\item MS-RPC DCOM buffer overflow
\item Ended up in the Blaster worm infecting : \url{https://en.wikipedia.org/wiki/Blaster_(computer_worm)}
\item notice the timeline, even if patches ARE available people didn't patch
\item This is why we have \emph{patch tuesdays}
\item Blaster was not as fast as SQL Slammer, which infected the internet in just minutes
\item Other example from book, NTPD, which is also a common and thought to be safe service to run
\item And \verb+mod_php+ nonterminating buffer vulnerability
\item Plus a few off-by-one errors in \verb+mod_rewrite+ and OpenBSD ftpd
\item ...
\end{list2}

\slide{Who do you trust?}

\begin{list2}
\item Example programs with flaws: OpenSSH, OpenSSL, Microsoft Windows, Linux
\item Who do you trust?
\item Can we trust any software?

\item What about Suricata? - CVE list\\
{\footnotesize\url{https://www.cvedetails.com/vulnerability-list/vendor_id-13364/product_id-27771/Openinfosecfoundation-Suricata.html}}
\item Spoiler: Bypass vulnerabilities, getting data past the IDS, Denial of Service crashing the IDS
\item Trust is also a design decision -- more on supply chain later
\end{list2}

\slide{Suricata IDS -- older CVEs}

\begin{list2}
\item CVE-2013-5919 Suricata < 1.4.6, CVSS 5.0: crash on malformed SSL record\\
      \url{https://www.cvedetails.com/cve/CVE-2013-5919/}
\item CVE-2014-6603 Suricata < 2.0.4, CVSS 5.0: crafted SSH banner, rule bypass,
      large allocation or OOB write\\
      \url{https://www.cvedetails.com/cve/CVE-2014-6603/}
\item CVE-2015-0971 Suricata < 2.0.8, CVSS 5.0: DER parser crash via
      SSL/TLS certificates\\
      \url{https://www.cvedetails.com/cve/CVE-2015-0971/}
\item CVE-2015-8954 Suricata < 2.0.6, CVSS 9.8: \verb+MemcmpLowercase+ skips
      first byte, IPS bypass via crafted HTTP\\
      \url{https://www.cvedetails.com/cve/CVE-2015-8954/}
\item CVE-2017-7177 Suricata < 3.2.1, CVSS 7.5: IPv4 defrag evasion,
      IP protocol not checked when matching fragments\\
      \url{https://www.cvedetails.com/cve/CVE-2017-7177/}
\item CVE-2017-15377 Suricata < 4.x, CVSS 7.5: redundant content checks
      until recursion limit (3000), DoS\\
      \url{https://www.cvedetails.com/cve/CVE-2017-15377/}
\end{list2}

\slide{Suricata IDS -- selected CVEs}

\begin{list2}
\item CVE-2018-6794 Suricata < 4.0.4, CVSS 5.3: TCP data before handshake,
      HTTP response signatures bypassed (Snort 2.9.4 too)\\
      \url{https://www.exploit-db.com/exploits/44247}
\item CVE-2024-32664 Suricata, CVSS 5.3: base64 out-of-bounds write
      via crafted traffic or datasets\\
      \url{https://github.com/OISF/suricata/security/advisories/GHSA-79vh-hpwq-3jh7}
\item CVE-2024-37151 Suricata, CVSS 7.5: fragments reusing IP ID break
      reassembly, policy bypass\\
      \url{https://github.com/OISF/suricata/security/advisories/GHSA-qrp7-g66m-px24}
\item CVE-2025-59147 Suricata, CVSS 7.5: crafted multiple SYNs, TCP
      detection bypass\\
      \url{https://bugzilla.redhat.com/show_bug.cgi?id=2400928}
\item CVE-2025-64330 Suricata, CVSS 7.5: heap over-read logging verdict
      in eve.alert (non-default config)\\
      \url{https://github.com/OISF/suricata/security/advisories/GHSA-83v7-gm34-f437}
\end{list2}

\begin{list1}
\item An IDS must parse hostile traffic -- and agree with the endpoint on how it should be understood
\end{list1}


\slide{Resources}

\begin{quote}\footnotesize
Microsoft Security Development Lifecycle (SDL) – The SDL is a software development process from Microsoft that helps developers build more secure software. It helps you address security compliance requirements while reducing development costs.

Open Web Application Security Project (OWASP) – OWASP is an online community that produces freely available articles, methodologies, documentation, tools, and technologies in the field of web application security.

Pushing Left, Like a Boss – A series of online articles that outlines different types of application security activities that developers should complete to create more secure code.

Microsoft identity platform – The Microsoft identity platform is an evolution of the Azure AD identity service and developer platform. It’s a full-featured platform that consists of an authentication service, open-source libraries, application registration and configuration, full developer documentation, code samples, and other developer content. The Microsoft identity platform supports industry-standard protocols like OAuth 2.0 and OpenID Connect.

Security best practices for Azure solutions – A collection of security best practices to use when you design, deploy, and manage cloud solutions by using Azure. This paper is intended to be a resource for IT pros. This might include designers, architects, developers, and testers who build and deploy secure Azure solutions.

Security and Compliance Blueprints on Azure – Azure Security and Compliance Blueprints are resources that can help you build and launch cloud-powered applications that comply with stringent regulations and standards.

Next steps\\
In the following articles, we recommend security controls and activities that can help you design, develop, and deploy secure applications.

* Design secure applications
* Develop secure applications
* Deploy secure applications
\end{quote}

Source:
\url{https://docs.microsoft.com/en-us/azure/security/develop/secure-dev-overview}

\slide{Architecture for modern-day web applications}

The average modern-day web application probably makes use of several of the follow‐
ing technologies:
\begin{list2}
\item  REST API
\item  JSON or XML
\item  JavaScript
\item  SPA framework (React, Vue, EmberJS, AngularJS)
\item  An authentication and authorization system
\item  One or more web servers (typically on a Linux server)
\item  One or more web server software packages (ExpressJS, Apache, NginX)
\item  One or more databases (MySQL, MongoDB, etc.)
\item  A local data store on the client (cookies, web storage, IndexDB)
\end{list2}
Source: \emph{Web Application Security}, Andrew Hoffman, 2020, ISBN: 9781492053118

Note: the list could have come from many places, very similar structure


\slide{Enterprise Integration Patterns}

\hlkimage{14cm}{eip-patterns.png}
Source: \url{https://www.enterpriseintegrationpatterns.com/}
Use patterns when building systems!


\slide{OSI Model and Internet Protocols}

\hlkimage{10cm,angle=90}{images/compare-osi-ip.pdf}



\slide{Document or it didn't happen}

\hlkimage{12cm}{mythical-man-month-documentary.png}

\begin{list2}
\item Make sure to document!
\item There are so many parts, so many decisions, soo many places thing can mess up
\item Have a minimum of documentation -- and the source files are not enough
\item Hint:  you can read the Mythical Man Month at \link{https://archive.org/details/MythicalManMonth}
\end{list2}

\slide{Why Have Formal Documents}

\hlkimage{10cm}{mythical-formal-documents.png}
Source: \emph{The Mythical Man-Month (Anniversary Edition)}
by Frederick P. Brooks Jr.

\begin{list2}
\item I use Git version control for my documentation purposes, some is in Zim Personal Wiki \link{https://zim-wiki.org/}
\item I also use an IPAM \link{https://spritelink.github.io/NIPAP/}
\end{list2}



\slide{Operating Systems}

Computers need software -- operating systems, and applications run on an OS:
\begin{list2}
\item OS is the base for your experiments
\item Must be managed, installed at least, updated and packages installed on top
\item Some applications need a regular OS, others are fine with containers
\item Some applications like Python can even use a virtual environment
\end{list2}

Select good operating systems, make sure to configure it -- \textbf{hardened/secure configuration}

\slide{First advice use the modern operating systems}

\begin{list1}
\item Newer versions of Microsoft Windows, Mac OS X and Linux
\begin{list2}
\item Buffer overflow protection
\item Stack protection, non-executable stack
\item Heap protection, non-executable heap
\item \emph{Randomization of parameters} stack gap m.v.
\end{list2}
\item Note: these still have errors and bugs, but are better than older versions
\item OpenBSD has shown the way in many cases\\ \link{http://www.openbsd.org/papers/}
\end{list1}

\vskip 1cm

\centerline{Always try to make life worse and more costly for attackers}



\slide{Virtualization}

Computers can run multiple operating systems, and OS distributions at the same time:
\begin{list2}
\item Full virtual machine virtualization
\item Containers -- including containers platforms like Kubernetes
\item Use it, and create many virtual machines/containers
\item Switch context easily, one project - one (or more) VMs
\item Create a new VM in seconds -- keep one template VM updated that you can clone
\item Install crap -- delete a VM, retry
\item Performance of VMs and containers on modern CPUs is great! Almost like having two physical computers!
\end{list2}

Maybe sometime use different physical hosts -- one for internal use, one for external use


\slide{Virtualization Overview}

\hlkimage{20cm}{containers.png}



\slide{Containers}

\hlkimage{11cm}{containers_only.png}

\begin{list2}
\item Containers package an application and all its dependencies into a single portable unit
\item Share the host OS kernel -- lighter and faster to start than full virtual machines
\item Docker is the dominant developer-facing tool; images are built from a \verb+Dockerfile+ and stored in registries
\item Images are not Docker-specific -- the industry standard is the {\bf OCI image format}, supported by all modern tools
\item {\bf Compose files} allow multi-container applications (e.g.\ app + database) to be defined and run together
\item Pre-packaged systems that require less administration -- {\bf only configuration}
\end{list2}


\slide{Container Standards and Runtimes}

\begin{list2}
\item The {\bf Open Container Initiative (OCI)} defines three open standards:\\
image-spec, runtime-spec, and distribution-spec \link{https://opencontainers.org}
\item Any OCI-compliant image can be built with one tool (e.g.\ Buildah, Docker) and run with another -- no vendor lock-in
\item {\bf runc} is the reference low-level runtime implementing the OCI runtime-spec; most runtimes use it under the hood
\item Kubernetes communicates with runtimes via the {\bf Container Runtime Interface (CRI)} -- a gRPC-based plugin API
\item {\bf containerd} and {\bf CRI-O} are the two main CRI-compliant runtimes; Docker's direct integration was removed in Kubernetes v1.24 \link{https://kubernetes.io/docs/setup/production-environment/container-runtimes/}
\item CRI-O is purpose-built for Kubernetes -- lightweight, OCI-compliant, and the default in Red Hat OpenShift \link{https://cri-o.io}
\end{list2}



\slide{Kubernetes}

\hlkimage{24cm}{kubernetes-anywhere.png}

\begin{list2}
\item Kubernetes (K8s) is an open-source system for automating deployment and management of containers at scale
\item Organises containers into \textit{Pods}, grouped by \textit{Deployments}, and exposed via \textit{Services}
\item Provides automatic scaling, self-healing, and rolling updates with zero downtime
\item Suited to large, complex microservice architectures -- significant operational overhead for small projects
\item Managed offerings (GKE, EKS, AKS) reduce setup burden but increase cloud dependency
\end{list2}

\slide{Kubernetes Security is more than just the core}

\hlkimage{12cm}{rice-container-security-1.png}
Source: \emph{Container Security}, (CS) Liz Rice, O'Reilly, Apr 2020

\begin{list2}
\item Attacks from inside, attacks inside containers
\item Some just use up resources, some attack others, some attack the system from the inside
\item \emph{Hacking Kubernetes} has a whole \emph{Chapter 3 Container Runtime Isolation}
\end{list2}

\slide{Kata Containers – containers in lightweight VMs}

\begin{list1}
\item Open source container runtime that runs each pod inside its own lightweight virtual machine
\item Each pod gets its own guest kernel instead of sharing the host kernel with all other containers
\item For Kubernetes and the user it still looks and behaves like normal containers
\item Uses a hypervisor like QEMU, Cloud Hypervisor or Firecracker to start the VM
\item An agent inside the VM, \emph{kata-agent}, starts and manages the containers on request from the runtime on the host
\item Container files are shared from the host into the VM with virtio-fs
\item Confidential Containers: the VM runs with encrypted memory (Intel TDX, AMD SEV-SNP), so even the host operator cannot look inside
\end{list1}


\slide{Kata Containers 2026 vulnerabilities – host escapes}

\begin{list2}
\item CVE-2026-44210 (8.8) – pod annotations inject arguments into virtiofsd, which serves the host root filesystem into the guest. Fixed in 3.31.0
\item CVE-2026-47243 (8.7) – guest root to host root escape via virtio-fs in runtime-rs
\item CVE-2026-50540 (9.6) – an unvalidated configuration path annotation gives code execution on the host. Fixed in 4.0.0
\item Common pattern: pod annotations that are enabled by default are an attack surface, because whoever can create a pod can influence the host side
\end{list2}

\slide{Kata Containers 2026 vulnerabilities – guest and CoCo}

\begin{list2}
\item CVE-2026-24834 – a container user can modify the guest VM filesystem and get root in the guest (Cloud Hypervisor). Fixed in 3.27.0
\item CVE-2026-41326 (8.2) – symlinks bypass the CopyFile policy, so the host can overwrite binaries inside the guest, even in CVMs. Fixed in 3.29.0
\item CVE-2026-64676 (5.7) – two mem-agent ttRPC methods skip the agent policy check. Fixed in 4.0.0
\item CVE-2026-77176 – genpolicy does not validate CreateContainer mount and storage rules well enough, so a malicious host operator can cause information disclosure
\item Advisories: \\ \url{https://github.com/kata-containers/kata-containers/security/advisories}
\end{list2}


\slide{Summary: Technology Comparison}

\begin{center}
\begin{tabular}{|l|l|l|l|}
\hline
{\bf Technology} & {\bf Best For} & {\bf Complexity} & {\bf Scalability} \\
\hline
VPS              & Single apps, full OS control   & Low    & Vertical (manual)   \\
Containers       & Portable, consistent deploys   & Medium & Horizontal (manual) \\
Kubernetes       & Large-scale microservices      & High   & Horizontal (auto)   \\
Virtual Hosts    & Multi-site on one server       & Low    & Limited             \\
\hline
\end{tabular}
\end{center}

\begin{list2}
\item Above does not consider if you are a private person running systems for yourself, friends, family
\item If you are a company with a few services
\item If you are a large company that want to switch over to something privately owned
\item Above are in my opinion good options for running your complete infrastructure!
\end{list2}


\slide{OWASP Kubernetes Security Cheat Sheet:\\
Container Sandboxing}

\begin{quote}
Cgroups and namespaces isolate containers, but the kernel still presents a large attack surface
\end{quote}

Recommends an extra sandbox layer for multi-tenant and highly untrusted clusters, with these open source options:
\begin{list2}
\item Kata Containers – stripped-down VMs with a small resource footprint
\item gVisor – its own kernel written in Go, sitting between the container and the host kernel
\item Firecracker – a very lightweight VM locked down with seccomp, cgroups and namespaces
\end{list2}

The OWASP Secure AI Model Ops Cheat Sheet also recommends microVMs, gVisor, Kata Containers or confidential compute for high-sensitivity workloads

Source: {\footnotesize\\
\emph{Kubernetes Security Cheat Sheet}, OWASP Cheat Sheet Series \\
\url{https://cheatsheetseries.owasp.org/cheatsheets/Kubernetes_Security_Cheat_Sheet.html}}


\slide{Virtual Hosts in Web Servers}

\begin{alltt}\footnotesize
server \{ server_name example.com www.example.com; root /var/www/example.com; \}
server \{ server_name example.org; root /var/www/example.org; \}
\end{alltt}


\begin{list2}
\item One server with one IP address can serve many websites
\item DNS for both {\bf example.com} and {\bf example.org} points to the same IP
\item The browser sends the name it wants in the HTTP {\bf Host} header, the web server picks the matching site
\item With HTTPS the name is also sent in TLS as {\bf SNI}, so the server can present the right certificate
\end{list2}

Side note: Digital Ocean and NGINX as a reverse proxy

That said DigitalOcean has very good documentation for lots of things! example:\\
\link{https://www.digitalocean.com/community/tutorials/how-to-configure-nginx-as-a-reverse-proxy-on-ubuntu-22-04}


\slide{NGINX as a Reverse Proxy}

\hlkimage{15cm}{nginx_reverse_proxy.png}
\begin{list2}
\item A reverse proxy sits in front of backend servers, accepting client requests and forwarding them internally
\item NGINX uses the \verb+proxy_pass+ directive to route traffic to upstream services by hostname or IP
\item Can provide TLS termination, load balancing, caching, and request buffering at the edge\\
A single system with NGINX get certificates with Let's Encrypt (or alternatives)
\end{list2}

% \link{https://en.wikipedia.org/wiki/Reverse_proxy}



\slide{Kubernetes Network parts -- NetworkPolicy}


\hlkimage{12cm}{cilium-overview.png}
\begin{list2}
\item Container Network Interface (CNI) and Cloud Native Computing Foundation (CNCF) project\\
\link{https://github.com/containernetworking/cni}
\item Multiple options -- popular and common ones Calico, Flannel and Cilium
\item Cilium - I have followed Cilium for a while, so I choose that
\end{list2}


\slide{Together with Firewalls - Virtual LAN (VLAN)}

\hlkimage{8cm}{vlan-portbased.pdf}

\begin{list1}
\item Managed switches often allow splitting into zones called virtual LANs
\item Most simple version is port based
\item Like putting ports 1-4 into one LAN and remaining in another LAN
\item Packets must traverse a router or firewall to cross between VLANs
\end{list1}

\slide{Virtual LAN (VLAN) IEEE 802.1q}

\hlkimage{16cm}{vlan-8021q.pdf}

\begin{list1}
\item Using IEEE 802.1q  VLAN tagging on Ethernet frames
\item Virtual LAN, to pass from one to another, must use a router/firewall
\item Allows separation/segmentation and protects traffic from many security issues
\end{list1}

\slide{Network Segmentation: Which VLANs to create}

Using VLANs to segment networks we can extend it to the rest of the network, and an example of isolated segments could be:
\begin{list2}
\item Guest network for cabled and Wi-Fi clients
\item Client networks can be isolated per department, floor or building
\item Server networks for multiple purposes – development, staging, testing, production, manufacturing, …
\item Internet servers in a demilitarized zone (DMZ) which would make it less likely that compromised internet
server would affect the rest of the network
\item Dedicated printer network
\item \textbf{Management network for virtualisation, one for network management, one for storage}
\end{list2}


\slide{Example, Restricting the Kubernetes API}

%\hlkimage{}{}

Watch out when learning, you may expose an \emph{admin interface}
\begin{quote}{\bf
By default, the Kubernetes API server listens on port 6443 on the first non-localhost network interface, protected by TLS}. In a typical product
ion Kubernetes cluster, the API serves on port 443. The port can be changed with the --secure-port, and the listening IP address with the --bind
-address flag.

...

If your cluster uses a {\bf private certificate authority, you need a copy of that CA certificate configured into your ~/.kube/config on the cli
ent}, so that you can trust the connection and be confident it was not intercepted.
\end{quote}

\begin{list2}
\item If an attacker can run Kubectl with credentials -- they can do anything
\item Defense in Depth is suggested -- so restrict at multiple levels
\item Only allow connections from specific network -- firewall filtering
\item Use TLS and verify/copy certificates
\end{list2}



\slide{Proxy servers and Web Application Firewalls (WAF)}

Filtering at higher layers is also possible
\begin{list2}
\item Web proxies for clients can help security a lot -- a centralized filter for everyone

\item Reverse proxies for web applications are called
Web Application Firewalls (WAF) -- and filter incoming web requests, and outgoing answers. Can help with attacks like SQL injection and exfiltration of data
\item Depending on your network it can replace or be combined with filtering on DNS servers, and I would prefer to filter domains with DNS
\item I would also prefer blocking large prefixes of IP destinations using routers/stateless packet filters -- maybe use BGP for distributing \emph{lists} or automation for downloading block lists
\end{list2}

Also great home lab projects!

\slide{CrowdSec -- Crowd sourced security information}

%\hlkimage{}{}

\begin{quote}
{\bf Detect and block Log4j exploitation attempts with CrowdSec}\\
If you work in Infosec, you had a very lousy weekend. And that’s because of the Log4j zero-day vulnerability (CVE-2021-44228) that was discovered. We had no cho
ice but to roll up our sleeves to help our community before things got messier than they already were.

As a result, we have released a scenario that will help you detect and block exploitation attempts of the vulnerability. This new scenario can be directly downl
oaded from our Hub and installed in a blink of an eye. Check this quick video to see the plugin in action:
\end{quote}

\url{https://www.crowdsec.net/blog/detect-block-log4j-exploitation-attempts}


\begin{list2}
\item Log4j is a popular software library in the Java world
\item CrowdSec quickly provided a list of systems scanning the internet for this vulnerability\\
 \link{https://gist.github.com/blotus/f87ed46718bfdc634c9081110d243166}
\item Full disclosure they gave me a hoodie at an event \smiley
\end{list2}


\slide{CrowdSec with NGINX}

%\hlkimage{8cm}{crowdsec_nginx.png}

\begin{list2}
\item {\bf CrowdSec} is an open-source, community-driven security engine -- the llama company \link{https://crowdsec.net}
\item {\bf NGINX Bouncer} -- a Lua plugin enforcing ban decisions in real time via \verb+access_by_lua_block+,
blocking malicious IPs before requests reach the backend
\item {\bf Security Engine} -- reads NGINX logs and detects behavioural patterns: HTTP probing,
scanning, brute force, and matches against the community blocklist
\item {\bf AppSec/WAF component} -- inline virtual patching, SQLi, XSS, and known CVE protection;
install with:
\verb+cscli collections install crowdsecurity/appsec-virtual-patching+
\item {\bf Community Blocklist} -- crowd-sourced IP reputation feed, continuously updated from
millions of sensors worldwide; free tier available
\item Two complementary modes: reactive (log-based, bans offenders) and
inline (WAF, blocks the request before it is processed)
\end{list2}
Source: Docs: \link{https://docs.crowdsec.net/u/bouncers/nginx/}


\slide{The OWASP Core Rule Set for WAF}

\begin{list2}
\item {\bf OWASP Core Rule Set (CRS)} -- the standard open source rules, protects against SQL injection, XSS, remote code execution and other common attacks \link{https://coreruleset.org}
\item CRS uses {\bf anomaly scoring} -- each matching rule adds points, the request is blocked when the total passes a threshold. {\bf Paranoia levels} 1--4 -- higher levels catch more attacks, but also block more legitimate traffic, start at level 1
\item Expect tuning: false positives are normal, run in detection-only mode first and read the logs
\item {\bf ModSecurity} -- open source Web Application Firewall engine, inspects each HTTP request and response against rules. Works with Apache, NGINX and IIS, version 3 (libmodsecurity) uses a connector module for NGINX
\item Trustwave ended support in July 2024, ModSecurity is now an OWASP project\\
\link{https://owasp.org/www-project-modsecurity/}
\item {\bf Coraza} -- newer OWASP WAF engine written in Go, first stable release 2021, fully compatible with CRS v4\\
\link{https://coraza.io}

\end{list2}



\slide{Image hygiene}

\begin{list2}
\item Small base images -- fewer packages, fewer CVEs, less for an attacker to use\\
e.g. slim, Alpine, distroless -- no shell is a feature
\item Multi-stage builds -- compilers and build tools stay out of the runtime image
\item Run as a named non-root user -- \verb+USER app+ in the Dockerfile
\item Pin base images by digest, not only by tag -- tags can move
\item No secrets in images -- layers keep history, \verb+docker history+ shows it
\item Scan images -- e.g. Syft and Grype, or Trivy
\item Rebuild regularly -- base images get security updates, your image does not
\end{list2}


\slide{Keeping Container Images secure}

\begin{list2}
\item \faWrench\ \verb+anchore+ open-source project that provides a centralized service for inspection, analysis, and certification of container images
\link{https://github.com/anchore/anchore-engine}\\
``As of 2023, Anchore Engine is no longer maintained. There will be no future versions released. Users are advised to use Syft and Grype.''
\item \faWrench\ \verb+Syft+ \link{https://github.com/anchore/syft} SBOM generation and \faWrench\ \verb+Grype+ \link{https://github.com/anchore/grype} vulnerability scanning -- alternatives to Trivy
\item Allow direct download from the internet into your cluster, may become a problem\\
-- use your own registry or a pull-through proxy
\item Malicious people are typosquatting popular containers!
\item Supply chain attacks in general are a problem
\end{list2}


\slide{Inventory and Control of Software Assets}

\begin{quote}
CIS Control 2:\\
Inventory and Control of Software Assets\\
Actively manage (inventory, track, and correct) all software on the network so that only authorized software is installed and can execute, and that all unauthorized and unmanaged software is found and prevented from installation or execution.
\end{quote}

\begin{list1}
\item What versions of software do we depend on?
\item What software needs to be phased out, upgraded?
\item SBOM -- Software Bill of Materials, a machine-readable list of components\\
formats CycloneDX and SPDX, generated by tools like syft and trivy
\item EU Cyber Resilience Act -- manufacturers must keep an SBOM
\end{list1}

Source: Center for Internet Security CIS Controls 7.1 CIS-Controls-Version-7-1.pdf


\slide{Continuous Vulnerability Management}

\begin{quote}
CIS Control 3:\\
Continuous Vulnerability Management\\
Continuously acquire, assess, and take action on new information in order to identify vulnerabilities, remediate, and minimize the window of opportunity for attackers.
\end{quote}

\begin{list1}
\item Acquire -- CVE feeds, vendor advisories, OSV, scanners in CI and registry
\item Assess -- not all critical CVEs are critical for you: reachable? exposed? exploited?\\
CISA Known Exploited Vulnerabilities (KEV) catalog helps prioritize
\item Act -- patch, upgrade, mitigate, or accept the risk with a documented reason
\item Continuously -- an image that was clean last month is not clean today
\end{list1}

Source: Center for Internet Security CIS Controls 7.1 CIS-Controls-Version-7-1.pdf


\slide{OSV -- vulnerabilities and malicious packages}

\begin{list2}
\item OSV -- Open Source Vulnerabilities, a database and a format, covering many ecosystems\\
npm, PyPI, Maven, Go, crates.io, Debian, ...
\item Vulnerabilities have IDs like \verb+CVE-+ and \verb+GHSA-+
\item Malicious packages have IDs starting with \verb+MAL-+\\
from the OpenSSF Malicious Packages repository
\item \verb+osv-scanner+ -- scans lockfiles, SBOMs and directories against OSV
\item \verb+npm audit+ -- uses the GitHub Advisory Database
\end{list2}

{\footnotesize\link{https://osv.dev/}\\
\link{https://github.com/ossf/malicious-packages}\\
\link{https://google.github.io/osv-scanner/}}


\slide{Overview virtualization}
\hlkimage{17cm}{zencurity-virtualisering.png}

Example servers with KVM and Mixed VM and Kubernetes

\slide{Overview services}
\hlkimage{16cm}{zencurity-services.png}

\slide{Kubernetes setup only VPN access to Nextcloud}
\hlkimage{13cm}{zencurity-k8s-netvaerk.png}
I do not use all Nextcloud features, so it is only reachable through WireGuard VPN

\slide{Real World: Upgrading Nextcloud on Kubernetes}

Critical PHP security advisory -- PHP 8.3 to 8.4, Nextcloud 33 to 34, about 30 minutes of work:
\begin{list2}
\item Behind a VPN there is no rush -- I could upgrade the next day at a good time, not late in the evening
\item Run applications in isolated rooms: own Kubernetes namespace, strict NetworkPolicy with both ingress and egress filters
\item Full backup first: a script dumping the 3 databases, and tar of all Nextcloud files
\item Check the kernel log with \verb+journalctl -k+ for PHP segfaults -- nothing found
\item Look for new PHP files from the last week -- not a guarantee, but cheap to run
\item Fetch the latest Helm chart for Nextcloud All-in-One, diff the settings for new required values
\item Start the upgrade and wait
\item The fulltextsearch app had a missing import, a one-line hotpatch\\
\link{https://github.com/nextcloud/fulltextsearch/issues/1005}
\item Wait a bit more, done
\end{list2}

\begin{alltt}\footnotesize
helm upgrade nextcloud-aio nextcloud-aio/nextcloud-aio-helm-chart --version 14.4.1
\end{alltt}

\slide{All Software Has Errors -- Know What You Run}

\begin{list2}
\item Every package, library and container image will get vulnerabilities -- also the ones you forgot you installed
\item Scanning tells you what to fix, patching does the fixing -- you need both
\item Remove what you do not use -- fewer packages, fewer findings
\item {\bf SBOM} -- Software Bill of Materials, a list of everything installed on a system or in an image
\item Example tools: {\bf Grype} matches packages against known vulnerabilities, {\bf Syft} builds SBOMs -- both open source from Anchore\\
\link{https://github.com/anchore/grype}
\item Debian hosts: the list of installed apt packages goes straight into Grype
\item Kubernetes: Syft builds an SBOM for each container image, then Grype scans the SBOM\\
The image is only downloaded once per version
\end{list2}

\slide{Example: Scanning My Own Infrastructure}

\begin{list2}
\item Debian and KVM hosts: a scanner VM collects package lists from 19 hosts over SSH,\\
using a restricted user that can only run a read-only export
\item Grype runs nightly, only High/Critical with a fix available
\item Push notification with ntfy, only when something new shows up
\item First run: 8367 findings -- after removing old leftover packages: 827
\item 37 Kubernetes images: 1274 findings, nginx:alpine went from 117 to 3 just by pulling the current image
\item Floating tags like \verb+latest+ bring fixes but change without review -- pinned versions are controlled but go stale
\item Do not take a new version the minute it is published -- most malicious versions are found within days
\end{list2}

\slide{Installed Is Not the Same as Running}

\begin{list2}
\item The scanner sees what is {\bf installed} -- a patched kernel package looks clean, but the old kernel is still {\bf running} until you reboot
\item Simple check: compare the running kernel with the newest installed one
\item If they differ, the host needs a reboot -- the nightly scan sends a notification
\item Same problem for libraries: a running service keeps using the old version until it is restarted, \verb+needrestart+ on Debian can show which
\item Kubernetes nodes: kured reboots nodes one at a time, inside a reboot window
\end{list2}

The dpkg log shows the install time, and the last boot shows when the new kernel actually started running:
\begin{alltt}\footnotesize
for h in cheese01 cheese02 gouda01; do
  echo "== $h"
  ssh $h 'zgrep -h " install linux-image-6.12.113" /var/log/dpkg.log*; uptime -s'
done
\end{alltt}

\breakdance


\slide{Session 1: Exercises: PMD Static Analysis}

\exercise{ex:pmd-static}


\slide{Session 2: Teaching: Sample Infrastructure, API Versioning and Containers}


\slide{Microservices for Java chapter 1}

\hlkimage{3cm}{microservices-for-java-developers.jpg}

Microservices for Java Developers , Christian Posta, 2016 O’Reilly\\
 \link{https://www.oreilly.com/programming/free/files/microservices-for-java-developers.pdf}

We will use the introduction, which is recommended reading.

\slide{Whats in the book}

\begin{quote}
This book is for Java developers and architects interested in developing microservices. We start the book with the high-level understanding and fundamental prerequisites that should be in place to be successful with a microservice architecture.
\end{quote}

Introducing some Java frameworks
\begin{list2}
\item The Spring ecosystem, Dropwizard and WildFly Swarm, we’ll use JBoss Forge CLI
\item Finally, when we build and deploy our microservices as Docker containers running inside of Kubernetes
\item They use VirtualBox with Docker and Kubernetes, YMMV
\item With source on Github, not updated in years though!\\
\link{https://github.com/redhat-developer/microservices-by-example-source}
\end{list2}
Source: {\footnotesize\\
Microservices for Java Developers , Christian Posta, 2016 O’Reilly}



\slide{Open source is also leading the charge in the technology space}

\begin{quote}\small
Open source is also leading the charge in the technology space. Following the commoditization curves, open source is a place developers can go to challenge proprietary vendors by building and
  innovating on software that was once only available (without source
  no less) with high license costs. This drives communities to build
  things like operating systems (Linux), programming languages (Go),
  message queues (Apache ActiveMQ), and web servers ( httpd ). Even
  companies that originally rejected open source are starting to come
  around by open sourcing their technologies and contributing to
  existing communities. As open source and open ecosystems have
  become the norm, we’re starting to see a lot of the innovation in
  software technology coming directly from open source communities
  (e.g., Apache Spark, Docker, and Kubernetes).
\end{quote}

\begin{list2}
\item We have used Linux and multiple products from the Apache website/foundation
\end{list2}
Source: {\footnotesize\\
Microservices for Java Developers , Christian Posta, 2016 O’Reilly}




\slide{Building distributed systems is hard}

\hlkimage{8cm}{microservices-java-book.png}

\begin{quote}\small
  Microservice architecture (MSA) is an approach to building software systems that decomposes business domain models into smaller,
  consistent, bounded-contexts implemented by services. These services are isolated and autonomous yet communicate to provide some
  piece of business functionality. Microservices are typically implemented and operated by small teams with enough autonomy that
  each team and service can change its internal implementation
  details (including replacing it outright!) with minimal impact across
  the rest of the system.
\end{quote}

Source: {\footnotesize\\
Microservices for Java Developers , Christian Posta, 2016 O’Reilly}





\slide{Teamwork}


\begin{list2}
\item Teams communicate through promises
\item specify these promises with interfaces
of their services and via wikis that document their services
\item Each team would be responsible for designing the service, picking
the right technology for the problem set, and deploying, managing
and waking up at 2 a.m. for any issues
\item Understand what the service is doing without being tangled into
other concerns in a larger application
\item Quickly build the service locally
\item Pick the right technology for the problem (lots of writes? lots of
queries? low latency? bursty?)
\item Test the service
\item Build/deploy/release at a cadence necessary for the business,
which may be independent of other services
\item Identify and horizontally scale parts of the architecture where
needed
\item Improve resiliency of the system as a whole



\end{list2}
Source: {\footnotesize\\
Microservices for Java Developers , Christian Posta, 2016 O’Reilly}




\slide{Challenges}

\begin{list2}
\item Microservices may not be efficient. It can be more resource intensive.
\item You may end up with
what looks like duplication.
\item Operational complexity is a lot higher.
\item It becomes very difficult to understand the system holistically.
\item It becomes significantly harder to debug problems.
\item In some areas you may have to relax the notion of transaction.
\end{list2}
Source: {\footnotesize\\
Microservices for Java Developers , Christian Posta, 2016 O’Reilly}




\slide{Design for Faults}

\begin{quote}
  Things will fail, so we must develop our applications to be resilient and handle failure, not just prevent it. We should
  be able to deal with faults gracefully and not let faults propagate to
  total failure of the system.

  Building distributed systems is different from building shared-
memory, single process, monolithic applications. One glaring differ‐
ence is that communication over a network is not the same as a local
call with shared memory.
\end{quote}

\begin{list2}
\item Networks are inherently unreliable
\end{list2}
Source: {\footnotesize\\
Microservices for Java Developers , Christian Posta, 2016 O’Reilly}


\slide{Design with Dependencies in Mind}

\begin{quote}
  To be able to move fast and be agile from an organization or
  distributed-systems standpoint, we have to design systems with
  dependency thinking in mind; we need loose coupling in our teams,
  in our technology, and our governance.
\end{quote}

Source: {\footnotesize\\
Microservices for Java Developers , Christian Posta, 2016 O’Reilly}



\slide{Design with Promises in Mind}

\begin{quote}
  In a microservice environment with autonomous teams and serv‐
  ices, it’s very important to keep in mind the relationship between
  service provider and service consumer. As an autonomous service
  team, you cannot place obligations on other teams and services
  because you do not own them; they’re autonomous by definition. All
  you can do is choose whether or not to accept their promises of
  functionality or behavior. As a provider of a service to others, all you
  can do is promise them a certain behavior.
\end{quote}

\begin{list2}
\item Promises as published by APIs and versions in those!
\item References \emph{Consumer-Driven Contracts: A Service Evolution Pattern}\\
\link{https://martinfowler.com/articles/consumerDrivenContracts.html}
\end{list2}
Source: {\footnotesize\\
Microservices for Java Developers , Christian Posta, 2016 O’Reilly}

\slide{Chapter 10: Service API and Contract Versioning\\
with Web Services and REST Services}

\begin{quote}
After a service contract is deployed, consumer programs will naturally begin forming dependencies on it. When we are subsequently forced to make changes
 to the contract, we need to figure out:
\begin{list2}
\item Whether the changes will negatively impact existing (and potentially future) service consumers
\item How changes that will and will not impact consumers should be implemented and communicated
\end{list2}
\end{quote}

Included in the chapter are references to:
\begin{list2}
\item Versioning Web Services
\item Versioning REST Services
\item Fine and Coarse-Grained Constraints
\end{list2}
Source: {\footnotesize\\
\emph{Service‑Oriented Architecture: Analysis and Design for Services and Microservices}, Thomas Erl, 2017}

\slide{Versioning and Compatibility}

\begin{list2}
\item Backwards Compatibility\\
A backwards-compatible change to a REST-compliant service contract might involve adding some new resources or adding new capabilities to existing resources. In each of these cases the existing service consumers will only invoke the old methods on the old resources, which continue to work as they previously did.
\item Forwards Compatibility\\
When a service contract is designed in such a manner so that it can support a range of future consumer programs, it is considered to have an extent of forwards compatibility. This means that the contract can essentially accommodate how consumer programs will evolve over time.
\end{list2}

Source: {\footnotesize\\
\emph{Service‑Oriented Architecture: Analysis and Design for Services and Microservices}, Thomas Erl, 2017}



\slide{Kubernetes API Versioning -- excerpt}

\begin{list1}
\item Alpha, e.g. \emph{v1alpha1}
\begin{list2}
\item Disabled by default, must be enabled explicitly
\item May contain bugs, can be dropped at any time without notice
\item May change incompatibly, only for short-lived test clusters
\end{list2}
\item Beta, e.g. \emph{v2beta3}
\begin{list2}
\item Disabled by default (except betas introduced before 1.22)
\item Well tested, the feature will not be dropped, but details may change
\item Incompatible changes come with migration instructions
\item Max lifetime: 9 months or 3 releases to deprecation, and the same again to removal
\end{list2}
\item Stable, e.g. \emph{v1}
\begin{list2}
\item Available in all future releases within a major version
\end{list2}
\end{list1}

Source: {\footnotesize\\
\emph{API Overview}, Kubernetes Documentation, CC BY 4.0
\url{https://kubernetes.io/docs/reference/using-api/}}



\slide{Docker }

%\hlkimage{}{}

\begin{quote}
  Package Software into Standardized Units for Development, Shipment and Deployment
  A container is a standard unit of software that packages up code and all its dependencies so the application runs quickly and reliably from one computing environment to another. A Docker container image is a lightweight, standalone, executable package of software that includes everything needed to run an application: code, runtime, system tools, system libraries and settings.
\end{quote}
Source: \\{\footnotesize
\url{https://www.docker.com/resources/what-container}}

\begin{list2}
  \item One of the most popular deployment methods today
\end{list2}

\slide{Containerized Applications}

\hlkimage{8cm}{container-what-is-container.png}

Source: {\footnotesize
\url{https://www.docker.com/resources/what-container}}

\begin{list2}
  \item See also \link{https://en.wikipedia.org/wiki/Linux_namespaces}\\
   \emph{Various container software use Linux namespaces in combination with cgroups to isolate their processes, including Docker[11] and LXC.}
\end{list2}

\slide{Docker Images and layers}

\hlkimage{10cm}{container-layers.jpg}

\begin{quote}
A Docker image is built up from a series of layers. Each layer represents an instruction in the image’s Dockerfile. Each layer except the very last one is read-only.
\end{quote}

Source: {\footnotesize
\link{https://docs.docker.com/storage/storagedriver/}}


\slide{Session 3: Exercises: SBOM Scanning and XML Library in Python}


\exercise{ex:syft-grype-juiceshop}
\exercise{ex-python-library}
\exercise{ex:whatweb-scanner}


\slidenext{}

\end{document}
